Five large photographs of the solar photosphere (adapted from
www.spaceweather.com) and the appropriate (and on-scale) Stonyhurst grids will
be given to the students (Figure 3A). The photographs given to the students,
will not be annotated. Each photograph is accompanied by the angle $B_0$ and
$P_0$ of the day of the observation. The student should calculate the rotation
of the Sun by measuring the coordinates of at least 3 well recognized sunspots
as they follow the rotation of the Sun.

% FIGURE NEEDED: Figure 3A -- five photographs of the solar photosphere
% (1, 3, 5, 7, 9 May 2013) and the two corresponding Stonyhurst grids
% (April 28 to May 06, May 07 to May 15); 2013_da_sol.pdf, p.~3 of the paper.

\begin{enumerate}
\item Correctly draw the axis of rotation of the Sun. This can be done by
  drawing a straight line at an angle of $P_0$ degrees anticlockwise from the
  vertical for each photograph (note: the photographs are given in equatorial
  coordinates.) \hfill(4 Points)
\item Choose the correct Stonyhurst grid and place it on each photograph, so
  that the solar axis on the grid coincides with the solar axis of the
  photograph. Estimate the heliocentric longitude, ($\ell_\odot$), and
  heliocentric latitude ($b_\odot$) with the help of the grid. Write down
  these coordinates in Table 1 for each of the 5
  photographs. \hfill(12 Points)

  \begin{center}
  \begin{tabular}{l cc cc cc}
  \toprule
  Date & \multicolumn{2}{c}{Sunspot S1}
    & \multicolumn{2}{c}{Sunspot S2} & \multicolumn{2}{c}{Sunspot S3}\\
  & $\ell_\odot$ & $b_\odot$ & $\ell_\odot$ & $b_\odot$
    & $\ell_\odot$ & $b_\odot$\\
  \midrule
  May 1 & 54    & 17.5 & 33  & 18.5 & 30.5 & 16\\
  May 3 & 27    & 18   & 6   & 19   & 4    & 16\\
  May 5 & 3     & 18.5 & $-19$ & 18.5 & $-23$ & 17\\
  May 7 & $-24$ & 18   & $-46$ & 18   & $-49$ & 16\\
  May 9 & $-49.5$ & 19 & $-72$ & 18   & $-77$ & 17\\
  \bottomrule
  \end{tabular}
  \end{center}
\item Construct the diagram $\Delta\ell_\odot/\Delta t$ for each sunspot.
  \hfill(2 for each graph)\quad(6 Points)

  % FIGURE NEEDED: three point diagrams Delta-l_S1, Delta-l_S2, Delta-l_S3
  % (longitude drift in degrees vs day, 0--10 on the x axis, 0--120 on the
  % y axis); 2013_da_sol.pdf, p.~4 of the paper.
\item Calculate the rotation ($\Delta\ell_\odot/\Delta t$) of the Sun for each
  sunspot:
  \[
  P_i = 8 * \frac{360}{\Delta\ell_{\odot i}}, \qquad
  P_1 = 27.8~days, \qquad
  P_2 = 27.4~days, \qquad
  P_3 = 26.8~days
  \]
  \hfill(3 for each graph)\quad(9 Points)
\item Calculate the average period of the Sun
  \begin{equation*}
  P_\odot = (P_1 + P_2 + P_3)/3 = 27.3~days \tag*{(1.5 Point)}
  \end{equation*}
\end{enumerate}
